\begin{abstract}
Radiologists annotating chest CT occasionally leave a visible lung nodule unmarked. Such omissions matter twice: when only part of the training data can be re-read, the re-reading budget must find them, and when a detector's alert threshold is calibrated on patients labelled negative, an unmarked nodule in that pool raises the threshold for everyone. We study both effects on LUNA16. A re-reading policy that favours candidates near the top of their scan and candidates on which two readers disagree reaches 79 of 92 training omissions in expectation at half the budget, more than three times the yield of re-reading in proportion to reader marks at the full budget. These additional labels, however, do not yield more detected omissions at a fixed alert rate. In a follow-up analysis that keeps every detector score fixed, removing only the 14 calibration patients with a nodule in the official consensus lowers the threshold for 280 of 283 detectors and reverses the sign of the estimated gap between the two training methods. The definition of the negative population therefore deserves the same scrutiny as the training labels.
\end{abstract}
